\newpage
\cleardoublepage
\chapter{Introduction}\label{chapter:introduction}

This work investigates how a wildlife species detector small enough for embedded hardware can be trained when real photographs of many target species are scarce. The deployment context is a consumer binocular that names the mammal a user is looking at, on the device itself and without a network connection. That context fixes two hard constraints: a per-frame latency budget on a mobile system-on-chip, and a class universe of $225$ mammal species rather than the few dozen coarse classes of a general detection benchmark. Two training strategies for such a detector are compared here, supervision from a much larger fine-tuned teacher model and direct fine-tuning of the same lightweight student model on the wildlife dataset. Of those $225$ classes, $50$ retain fewer than $150$ usable real photographs each. This work therefore also asks whether diffusion-generated imagery can stand in for the missing real data, and how far a detector trained on it still transfers to real photographs.

%%%%%%%%%%%%%%%%%%%%%%%%%%%%%%%%%%%%%%%%%%%%%%%%%
\section{Motivation}\label{sec:motivation}

Species identification in the field has begun to move from printed field guides into the optics themselves. The \textit{AX Visio} binocular \cite{AXVisio} runs a species-recognition model on board and names the animal the user is already looking at. Naming has to happen while that animal is still in the field of view, which rules out sending the frame to a remote service. The habitats where the device is used rule that out a second time, since network coverage there cannot be assumed. The recognition model therefore has to fit a mobile system-on-chip and meet a per-frame latency budget on it.

The deployment target is the \textit{Qualcomm QCS605} \cite{QCS605}, a mobile platform combining a Hexagon 685 digital signal processor with an Adreno 615 graphics processor. Its compute budget excludes the model scales at which detection accuracy is normally reported. Development and benchmarking in this work targeted a \textit{Raspberry Pi 400} \cite{ltdBuyRaspberryPi} instead of the QCS605 itself, for the software-maturity reasons set out in \Cref{sec:lit_embedded_inference}.

Accuracy in object detection has largely been bought with scale, which is what makes an embedded target awkward. The teacher model used in this work carries roughly $196\,\mathrm{M}$ parameters, against $2.71\,\mathrm{M}$ for the student model actually exported to the device. Knowledge distillation is the established answer to a gap of that size. \citeauthor{buciluaModelCompression2006} \cite{buciluaModelCompression2006} trained a small, fast model to reproduce the outputs of a large ensemble, and \citeauthor{hintonDistillingKnowledgeNeural2015} \cite{hintonDistillingKnowledgeNeural2015} generalized the idea by having the student match the teacher's softened output distribution rather than its hard predictions. Whether that mechanism survives a move from general object classes to fine-grained species is the question this work was built around.

Two properties of the wildlife domain make that question non-trivial. The class universe is fine-grained, with several of the $225$ classes separable only by markings a detector has to learn to attend to \cite{weiFineGrainedImageAnalysis2022}. And the available imagery is long-tailed, with a few well-photographed species dominating the corpus while most species are supported by comparatively few usable photographs \cite{zhangDeepLongTailedLearning2023}. Machine-learning work on wildlife imagery treats both properties as recurring obstacles rather than as edge cases \cite{tuiaPerspectivesMachineLearning2022}. Camera-trap detection is also known to degrade sharply once a model is evaluated at locations it was not trained on \cite{beeryRecognitionTerraIncognita2018}, so a thin per-class image pool is a generalization problem and not only a sample-count problem.

Text-to-image diffusion models \cite{rombachHighResolutionImageSynthesis2022} offer the one intervention that can add appearance content a data-poor class's own photographs do not contain, such as an unrecorded pose, habitat or lighting condition. That offer comes at a measured discount. Training a classifier from scratch on generated images required roughly five times as many of them to match real images of the same classes \cite{heSyntheticDataGenerative2023}, and prompting an off-the-shelf generator with bare class names produced a very poor classifier \cite{aziziSyntheticDataDiffusion2023}. Earlier work by the present author found that generated image context conferred no downstream detection advantage over per-pixel random noise on a small-object car-detection task \cite{hedrichGeneratingSyntheticDatasets2026}. Synthetic imagery is therefore treated here as a supplement whose usefulness is measured rather than assumed.

Wildlife species detection for a handheld optic exercises both problems at once, which is what makes it a demanding test case. It combines a fixed latency ceiling, a fine-grained class universe, a long-tailed image supply, and a product requirement that the resulting weights remain commercially licensable. None of these constraints was constructed for the experiment. Each of them arrived with the use case.

%%%%%%%%%%%%%%%%%%%%%%%%%%%%%%%%%%%%%%%%%%%%%%%%%
\section{Objective}\label{sec:objective}

The objective of this work is to build, train and evaluate a nano-scale wildlife species detector for embedded deployment, and to measure which of the available training strategies actually helps it. Detection accuracy is reported throughout on the \textbf{mixed test set}, the union of the real and synthetic held-out images, with the \textbf{real test set} breakout stated alongside it every time (\Cref{sec:mixed_default}). Four research questions organize the work:

\begin{enumerate}
    \item Does knowledge distillation from a fine-tuned teacher model into a nano-scale student model yield higher detection accuracy than fine-tuning the same student directly on the wildlife dataset, given the domain shift from COCO-style classes to $225$ mammal species?
    \item Which synthetic-image generator produces the most useful training data for this domain, and do the two criteria available before any detector is trained, apparent visual realism and per-image price, predict that ranking?
    \item How does the composition of a class's training data, from synthetic images alone through a large pool of real photographs, affect detection accuracy, and does accuracy measured on the mixed test set hold up on the real test set alone?
    \item Does a nano-scale detector meet the per-frame latency budget of the target hardware without quantization or any other compression step, measured on a target-adjacent proxy device?
\end{enumerate}

The four questions did not receive equal effort, and the imbalance is reported here rather than smoothed over. Roughly $80\%$ of the available time went into dataset construction and synthetic-image generation, so questions $2$ and $3$ carry the bulk of the experimental work. Question $1$ is answered from a single teacher and student pairing at one distillation hyperparameter setting. Question $4$ is answered from a single full-precision runtime measurement on the proxy device, with no quantization pass behind it. \Cref{sec:results_kd_and_embedded} restates that reduced scope where those results are reported, and \Cref{sec:limitations} treats it as a named limitation of the work.

Each question is restated and answered in turn in \Cref{chapter:discussion_conclusion}, from the evidence assembled in \Cref{chapter:results}.

%%%%%%%%%%%%%%%%%%%%%%%%%%%%%%%%%%%%%%%%%%%%%%%%%
\section{Environment}\label{sec:environment}
This Master's thesis is carried out in partnership with \textit{inovex GmbH} \cite{inovex}, an IT consultancy headquartered in Karlsruhe, Germany, that positions itself at the forefront of digital transformation. \textit{inovex} maintains dedicated practice areas spanning \textit{Application Development}, \textit{IT Operations}, and \textit{Data Management \& Analytics}. The present work is hosted within the latter, which builds business intelligence, big-data, and applied AI solutions for its clients. The collaboration grants access to \textit{inovex}'s internal compute infrastructure, whose GPU resources make the training and distillation experiments in this work feasible within the available timeframe. Beyond providing this infrastructure, \textit{inovex} placed no fixed requirements on the research direction, leaving the choice and design of the technical approach entirely to the author.

\paragraph{Utilization of AI Tools in Thesis Preparation}
Several AI tools supported the preparation of this work, in line with the guidelines issued by \textit{Hochschule Karlsruhe} on May 16, 2024. ChatGPT \cite{ChatGPT}, DeepL Translator \cite{DeepLTranslateWorlds}, and DeepL Write \cite{DeepLWriteAIpowered} were used to polish sentence structure, correct spelling and grammar, and translate text the author had already written. Because this use only reformulates content the author originally authored, it falls outside the guideline's citation requirement for AI-generated content, but is disclosed here regardless for transparency.

Technical work in this work drew on additional AI assistance. \textit{Claude} \cite{Claude} and \textit{Claude Code} \cite{ClaudeCode}, Anthropic's models and coding agent, were used throughout the data-engineering, training, and evaluation pipelines to draft, review, and debug code under the author's direction. \textit{Gemini} \cite{GeminiAPI} supported occasional coding and research queries.

Three roles of that assistance are worth separating, because they carry different weights. The largest by volume was code generation. Most of the data-engineering and training code behind this work was drafted by \textit{Claude Code} against plans and design decisions written by the author, then reviewed, corrected and run by the author. This covers the curation and quality-filtering pipeline of \Cref{sec:data_quality_and_curation}, the synthetic-generation and generator-comparison tooling of \Cref{sec:synthetic_data_supplementation,sec:generator_comparison}, the training pipelines of \Cref{sec:model_training_experiment_tracking}, and the evaluation suite of \Cref{sec:evaluation_framework}.

The second role was iterative debugging and analysis. Defects were diagnosed in dialogue with the tool, and intermediate results were summarized and cross-checked with it. One consequence is reported in this work rather than hidden, namely the Band A data-loading defect of \Cref{sec:band_structure}, which was found and repaired through exactly that process and which invalidated a first generation of trained models.

The third role was prose drafting. Sections of this manuscript were drafted with \textit{Claude Code} from author-supplied source material, structure and argument, and then edited by the author. Drafting followed a fixed set of house-style rules covering sentence construction, terminology, claim-to-evidence traceability and citation discipline. Every citation was checked against the bibliography, and every reported number against the artifact that produced it.

As with the tools named above, all AI-assisted output was reviewed, tested, and corrected where necessary by the author. This work's research questions, experimental design, and interpretation of results remain the author's own. Where a conclusion is not supported by the evidence, this work states that in its own voice. No AI tool was treated as a source of factual claims about the literature, and no cited work is cited on an AI tool's account of it.

%%%%%%%%%%%%%%%%%%%%%%%%%%%%%%%%%%%%%%%%%%%%%%%%%
\subsection{Pre-Work and Preliminary Resources}\label{sec:pre_work}

Several resources and constraints were already fixed before this work's own experimental work began. They are recorded here because they bound the design space the later chapters operate in.

\paragraph{The Production Baseline}
The \textit{AX Visio} already ships a two-stage recognition pipeline, a fine-tuned \textit{YOLOv5s} detector followed by the \textit{SpeciesNet} \cite{gadotCropNotCrop2024a} classifier applied to each detected crop. The full \textit{MegaDetector} \cite{beeryEfficientPipelineCamera2019} was too large for the device and was replaced there by the smaller detector. That arrangement supplied both the accuracy reference this work measures against and the architecture family its student model was drawn from (\Cref{sec:lit_yolo_vs_ensemble,sec:model_selection}).

\paragraph{Class Universe and Species Selection}
The class universe was inherited rather than designed. It originates in \textit{SpeciesNet}'s published taxonomy, restricted to non-bird mammals and then filtered down by the AX Visio product owner to the labels a binocular user is plausibly able to observe (\Cref{sec:target_class_definition}). Inclusion decisions drew on observation counts scraped from \textit{GBIF} \cite{GBIF}, on the premise that a species with very few recorded observations is both hard to train on and unlikely to be encountered.

\paragraph{Licensing Constraint}
One constraint was settled before any model was trained. \textit{YOLOv5} \cite{ultralyticsComprehensiveGuideUltralytics} is commercially usable only up to commit \texttt{5cdad89}, after which Ultralytics changed its licensing terms. Because the resulting weights are intended for a commercial product, this bounded which codebase revision was usable at all, independently of any accuracy argument (\Cref{sec:model_selection}).

\paragraph{Image Sources Assessed in Advance}
Candidate corpora were surveyed before sourcing began. The \textit{iNaturalist} competition datasets \cite{hornINaturalistSpeciesClassification2018} were preferred as an open collection of daylight photographs with heavy overlap onto the target class list. Camera-trap archives were examined and largely set aside, since their imagery is shaped by a fixed trap rather than by a handheld optic (\Cref{sec:data_sourcing_and_taxonomy}).

\paragraph{Earlier Work on Synthetic Training Data}
Earlier work by the present author \cite{hedrichGeneratingSyntheticDatasets2026} pasted real car crops onto generated backgrounds to improve small-object detection, and found that a per-pixel random-noise control matched or exceeded the generated context at every substitution ratio tested. That result is what leaves generator choice an open question here rather than a settled one (\Cref{sec:generator_comparison}).

%%%%%%%%%%%%%%%%%%%%%%%%%%%%%%%%%%%%%%%%%%%%%%%%%
\section{Structure}\label{sec:structure}

The remainder of this thesis is organized into four chapters and an appendix.

The \hyperref[chapter:literature_review]{\textit{literature review}} surveys the work this work's technical decisions rest on, weighted to match where the effort actually went. It traces object detection from its two architectural lineages to the nano-scale YOLO variants relevant here, and closes that section with the constraint-driven choice of student architecture. Synthetic training data then receives the greatest depth, covering generative image models, the sim-to-real gap, synthetic-image quality evaluation, and the long-tail literature. Knowledge distillation and embedded inference feasibility follow at foundational depth, matching the reduced scope of the corresponding experiments. The chapter ends with the evaluation literature this work's own metrics draw on.

\hyperref[chapter:methods_implementation]{\textit{Methods and Implementation}} describes how the dataset and the experiments were built. Four sections cover the dataset itself, from class taxonomy and image sourcing through curation and ground-truth verification, synthetic supplementation for the data-poor classes, and training-time augmentation treated as a controlled variable. A self-contained sub-study then revisits the choice of image generator through a controlled comparison across ten generators and three prompt regimes. The chapter closes with the training and experiment-tracking setup, and with the evaluation framework that defines the axes every reported number is indexed by.

The \hyperref[chapter:results]{\textit{results}} chapter reports the measurements in the order of their weight in the work. It opens with the practicability of synthetic training data, judged through the generator comparison. It then reports detection performance across the training bands and the test domains, which is the dataset's own central comparison. A smaller closing section reports the single distillation-versus-direct-fine-tuning comparison and the runtime benchmark on the proxy device, explicitly scoped as basic rather than exhaustive. A final section collects cross-cutting observations that belong to no single experiment.

\hyperref[chapter:discussion_conclusion]{\textit{Discussion and Conclusion}} restates the four research questions of \Cref{sec:objective} and answers them one by one from the evidence of the results chapter. It then states the limitations of the work, proposes further work motivated by those limitations, and closes with a summary of the contribution.

The \hyperref[appendix]{\textit{appendix}} holds supplementary material kept out of the main chapters, including the per-class band assignment, the frozen look-alike group listing, and qualitative sample galleries showing each generator's output for the same species beside a real photograph.
